% Synthetic evaluation fixture, not a research submission.
\documentclass{article}
\usepackage{microtype,graphicx,booktabs,hyperref}
\usepackage{icml2026}
\usepackage{algorithm,algorithmic,amsmath,amssymb}
\newtheorem{PWTheorem}{Theorem}
\newcommand{\method}{StepScale}
\newcommand{\states}{M_{\mathrm{state}}}
\providecommand{\theHalgorithm}{\arabic{algorithm}}
\icmltitlerunning{StepScale: Optimizer-State Storage}
\begin{document}
\twocolumn[
\icmltitle{StepScale: An Optimizer-State Storage Evaluation}
\begin{icmlauthorlist}
\icmlauthor{Example Researcher}{example}
\end{icmlauthorlist}
\icmlaffiliation{example}{Synthetic Evaluation Lab, Example City, Example Country}
\icmlcorrespondingauthor{Example Researcher}{fixture@example.org}
\icmlkeywords{optimizer state, storage, evaluation fixture}
\vskip 0.3in]
\printAffiliationsAndNotice{}
\begin{abstract}
Optimizer state can limit the memory available for model training.
We introduce \method{}, which shares one scale per block to reduce stored optimizer state.
Across two tasks with matched training steps and hardware, \method{} reduces optimizer-state memory from 8 to 6 GB while matching the reported mean accuracies.
Mean training time increases from 100 to 105 seconds.
\end{abstract}
\section{Introduction}\label{sec:introduction}
Optimizer-state storage competes with the rest of the training workload for device memory.
\method{} shares one scale within each block to reduce this storage cost.
Our evaluation measures the resulting memory and training-time tradeoff on two tasks.

\section{Method and Analysis}\label{sec:method}
The supplied design shares one scale per block. We report optimizer-state memory separately from training time and accuracy.
For positive baseline memory $M_b>0$ and method memory $M_s$, define the fractional state-memory reduction as
\begin{equation}\label{eq:memory}
 r=\frac{M_b-M_s}{M_b},\qquad \states=M_s.
\end{equation}
The measurement protocol fixes the hardware and number of training steps for the two methods.

\begin{algorithm}[h]
\caption{Computing the Reported Storage Reduction}\label{alg:summary}
\begin{algorithmic}[1]
\REQUIRE Supplied state-memory values $M_b>0$ and $M_s$
\STATE Compute $r=(M_b-M_s)/M_b$.
\ENSURE Fractional reduction $r$
\end{algorithmic}
\end{algorithm}


\paragraph{Measurement Objects.}
State memory is the storage associated with the optimizer state.
The comparison uses this object consistently for both methods: the baseline uses 8 GB and \method{} uses 6 GB.
Training time measures elapsed time for the fixed step budget on the same hardware.
Accuracy is reported separately for Task A and Task B rather than combined across tasks.
These three measurements describe the storage, time, and predictive-quality dimensions of the same comparison.

\paragraph{Reporting the Tradeoff.}
Equation~\ref{eq:memory} uses baseline memory as its denominator.
A reduction from 8 to 6 GB therefore releases 2 GB, or one quarter of the baseline optimizer-state storage.
The training-time difference uses the baseline time in the same way: the increase from 100 to 105 seconds is 5 seconds, or 5\%.
Presenting the two changes separately makes the storage benefit and time cost directly comparable without assigning an arbitrary aggregate efficiency score.

\section{Results}\label{sec:results}
Each task uses five seeds. Table~\ref{tab:wide} lists the supplied means.
Under this setup, it may perhaps be possible that \method{} could potentially reduce optimizer-state memory from 8 to 6 GB, although this should be interpreted cautiously.
Both methods achieve mean accuracies of 80\% on Task A and 78\% on Task B.
Mean training time increases from 100 to 105 seconds.
The measurements support lower state storage at the reported mean accuracy levels.

\begin{figure}[h]
\centering
\fbox{\parbox{0.32\linewidth}{\centering Baseline\\8 GB}}
$\longrightarrow$
\fbox{\parbox{0.32\linewidth}{\centering \method{}\\6 GB}}
\caption{Supplied optimizer-state memory for the comparison. The boxes illustrate the two values.}\label{fig:storage}
\end{figure}

\begin{table*}[t]
\centering
\caption{Synthetic fixture values. Accuracy and time are supplied means from five seeds; memory measures optimizer-state storage.}\label{tab:wide}
\begin{tabular}{llrrr}
\toprule
Task & Method & Accuracy (\%) & State memory (GB) & Training time (s)\\
\midrule
A & Baseline & 80 & 8 & 100\\
A & \method{} & 80 & 6 & 105\\
B & Baseline & 78 & 8 & 100\\
B & \method{} & 78 & 6 & 105\\
\bottomrule
\end{tabular}
\end{table*}


The two tasks provide separate checks of the reported accuracy means.
On Task A, both methods reach a mean accuracy of 80\%; on Task B, both reach 78\%.
The storage reduction is the same in both comparisons.
The repeated accuracy means and the common storage values make the memory-time tradeoff consistent across these supplied task summaries.
A reader deciding between the two configurations can use the 2 GB reduction and the 5 second increase according to the resource constraint of the training job.

The comparison also illustrates why a single label such as ``faster'' would obscure the result.
The measured benefit is smaller optimizer-state storage.
Elapsed training time is larger, so the result should retain both measurements when summarizing the method.

\section{Limitations and Conclusion}
The evaluation covers two tasks on the tested hardware; distributed training remains untested.
\method{} reduces state storage in this comparison, with an increase in training time.
\label{pw:main-end}

\section*{Impact Statement}
This synthetic fixture tests writing and typesetting behavior; the values are supplied examples rather than measurements from a deployed system.
\bibliographystyle{icml2026}\bibliography{references}
\clearpage\appendix
\section{Supporting Details}\label{app:supporting}
The example values give a fractional storage reduction of $r=(8-6)/8=0.25$ using Equation~\ref{eq:memory}.
Algorithm~\ref{alg:summary} states the arithmetic used to report the reduction; it is not an optimizer update rule.
Figure~\ref{fig:storage} and Table~\ref{tab:wide} report the same supplied state-memory values.

\subsection{Mathematical Typesetting Check}
The following identity provides a multiline-equation and theorem check.
\begin{equation}\label{eq:identity}
\begin{aligned}
\sum_{i=1}^{n}(x_i-\bar{x})^2
 &=\sum_{i=1}^{n}x_i^2-2\bar{x}\sum_{i=1}^{n}x_i+n\bar{x}^2\\
 &=\sum_{i=1}^{n}x_i^2-n\bar{x}^2,
\end{aligned}
\end{equation}
where $n\geq1$, $x_i\in\mathbb{R}$, and $\bar{x}=n^{-1}\sum_{i=1}^{n}x_i$.
\begin{PWTheorem}[Nonnegative Centered Sum]
For $n\geq1$ and real $x_1,\ldots,x_n$, the left-hand side of Equation~\ref{eq:identity} is nonnegative.
\end{PWTheorem}
\textit{Proof.} Each squared real term is nonnegative, so their sum is nonnegative.

\subsection{Fixture Provenance}
The official author guide~\cite{formatguide} is cited to check bibliography rendering, not to support the synthetic research values.

\end{document}
